\section{Задача Коллатца: двигатель на ограниченном топливе}\label{sec:collatz}

Все прежние маски читали трудную задачу через \emph{суррогатный} двигатель. Коллатц --- исключение:
здесь двигатель есть \emph{буквальный} объект на $\N$. Мы переводим гипотезу на язык двигателя и
каната, находим, где именно она нарушает закон спуска, и доводим зелёную механику до именованного
закона каната --- который в универсальной форме мы здесь же машинно \emph{опровергаем}. Это
опровержение не случайно: закон был принят четвёртой границей первопричины \axiom{}, и его падение ---
центральное событие раздела.

\subsection{Карта и где ломается закон спуска}

\begin{definition}[\code{CollatzMap}]\label{def:collatzMap}
Ускоренная карта Коллатца $T\colon\N\to\N$ вкладывает вынужденное деление на 2 после нечётного шага в
себя:
\begin{equation}\label{eq:collatz-map}
T(n)=\begin{cases}n/2,& n\equiv 0\ (\mathrm{mod}\ 2),\\[2pt] (3n+1)/2,& n\equiv 1\ (\mathrm{mod}\ 2).\end{cases}
\end{equation}
Чётный шаг --- спуск ровно на два (краевой случай $A=2$, не строгий); нечётный шаг --- \emph{впрыск
топлива}: «$+1$» в $3n+1$ поднимает высоту.
\end{definition}

Двигатель Евклида останавливается гарантированно, ибо каждый шаг --- строгий спуск, откуда
\code{no\_infinite\_descent} вынуждает $H(t)<H_0/A^{t}<1$. У Коллатца такого закона нет.

\begin{theorem}[\code{odd\_accel\_ascends}]\label{thm:odd_accel_ascends}
Для нечётного $n\ge 1$ ускоренный шаг строго поднимает высоту: $n<\tfrac{3n+1}{2}=T(n)$. \green{}
\end{theorem}

\begin{theorem}[\code{not\_descent\_on\_odd}]\label{thm:not_descent_on_odd}
Для нечётного $n\ge 1$ шаг не является спуском, $\neg\,(T(n)<n)$. Значит, высота вдоль орбиты не
\code{StrictAnti}, и \code{no\_infinite\_descent} неприменим. \green{}
\end{theorem}

\emph{Идея доказательства.} Оба факта --- однострочная арифметика на $\N$. Численно высота гуляет вверх
на четыре-пять порядков до спуска (пик/старт $\approx 5{,}4\times 10^{4}$ при $n=159487$); строгого
монотонного Ляпунова нет. Это точная граница между доказанной остановкой двигателя Евклида и
открытостью Коллатца: впрыск топлива временно поднимает высоту, поэтому строгая форма второго закона
сломана.

\begin{theorem}[\code{no\_monotone\_height}]\label{thm:no_monotone_height}
Восходящие позиции неограниченны: для всякого $N$ найдётся $n>N$ с $n\equiv 1\ (\mathrm{mod}\ 4)$ и
$n<(3n+1)/2$ (свидетель $n=4N+5$). Поэтому монотонной $\log_2$-высоты нет. \green{}
\end{theorem}

В среднем топливо нетто-расходуется: средний геометрический дрейф за шаг равен $\sqrt{3}/2\approx
0{,}864<1$, а баланс halvings/triplings $=2{,}016$ превышает порог $\log_2 3=1{,}585$. Но это закон
\emph{среднего}, а гипотеза требует, чтобы останавливалась \emph{каждая} траектория. Наша
высокоэнергетическая несовместимость требует строгого спуска на каждом шаге; Коллатц --- двигатель на
ограниченном топливе по энергии, но не EPMI-двигатель.

\subsection{Перетягивание каната: мультипликативный бюджет окна}

Читаем динамику без логарифмов. Двигатель продвигается не более чем на один двоичный ранг за шаг;
канат тянет назад ровно на один ранг (деление пополам) или на два ($n\equiv 0\ \mathrm{mod}\ 4$). Тогда
спуск --- бухгалтерское неравенство между числом чётных и нечётных шагов в окне.

\begin{theorem}[\code{engine\_at\_most\_one\_rank}]\label{thm:engine_at_most_one_rank}
Один ход двигателя поднимает не более чем на ранг: $T(n)\le 2n$ для всех $n$. \green{}
\end{theorem}

\begin{theorem}[\code{window\_budget}]\label{thm:window_budget}
За окно из $k$ шагов, среди которых $t=\mathrm{oddCount}(k,n)$ нечётных и
$e=\mathrm{evenCount}(k,n)$ чётных, пошаговые границы перемножаются при любом порядке шагов:
\begin{equation}\label{eq:window-budget}
2^{e}\cdot T^{k}(n)\;\le\;2^{t}\cdot n.
\end{equation}
Это чистая $\N$-арифметика: floor-деления съедены точностью пошаговых границ, логарифмы не нужны. \green{}
\end{theorem}

\begin{theorem}[\code{window\_descends}]\label{thm:window_descends}
Если канат перетянул окно строго ($t<e$) и $n\ge 1$, то $T^{k}(n)<n$. \green{}
\end{theorem}

\begin{proof}
Из \cref{eq:window-budget}: при $t<e$ имеем $2^{t+1}\le 2^{e}$, откуда $2\cdot T^{k}(n)\le n$ и
$T^{k}(n)<n$.
\end{proof}

Взят точный порог $e>t$ --- а не неправдоподобный в среднем $e>2t$ --- ровно тот, что зелёная граница
$\times 2$ конвертирует в спуск. Отказ от $\log_2 3$ стоит запаса $\sim 1{,}6\%$ (настоящий спуск
требует лишь $e>0{,}585\,t$, но это уже логарифмы, не core Lean).

\begin{definition}[\code{RopeCountingLaw}]\label{def:ropeCountingLaw}
Для $n$ закон доминирования каната гласит, что над поглощающим циклом каждая позиция открывает окно с
чётным перевесом:
\begin{equation}\label{eq:rope-law}
\mathrm{RopeCountingLaw}(n)\ :\equiv\ \forall j,\ \bigl(2<T^{j}(n)\Rightarrow \exists k,\ \mathrm{oddCount}(k,T^{j}(n))<\mathrm{evenCount}(k,T^{j}(n))\bigr).
\end{equation}
\end{definition}

\begin{theorem}[\code{reaches\_one\_of\_countingLaw}]\label{thm:reaches_one_of_countingLaw}
Если $n\ge 1$ и выполнен $\mathrm{RopeCountingLaw}(n)$, то $\exists K,\ T^{K}(n)=1$. \green{}
\end{theorem}

\emph{Идея доказательства.} Лестница целиком зелёная: $\mathrm{RopeCountingLaw}(n)$ питает бюджет окна
(\cref{thm:window_budget}), тот даёт закон ценностного спуска, а он индукцией по потолку значения
доводит до $1$. Это \emph{условная} теорема о конкретном $n$ --- «дай закон для этого $n$, и получишь
остановку». Что теперь нельзя --- раздать посылку всем $n$ сразу; универсальная форма опровергнута ниже.

\subsection{Кованое опровержение: карабкающаяся орбита $27$}

Закон каната требует чётного перевеса в \emph{каждой} позиции над циклом, включая самый старт $j=0$.
Орбита $27$ --- самое знаменитое число в фольклоре Коллатца --- карабкается до пика $4616$, прежде чем
сдаться, и такого окна из старта не открывает ни разу: пока число лезет вверх, шаги нечётно-тяжёлые, и
разность $\mathrm{evenCount}-\mathrm{oddCount}$ в префиксных окнах никогда не становится положительной.

\begin{theorem}[\code{counts\_at\_70\_at\_27}]\label{thm:counts_at_70_at_27}
Орбита $27$ приходит к единице за $70$ ускоренных шагов:
\begin{equation}\label{eq:counts-27}
\mathrm{oddCount}(70,27)=41,\qquad \mathrm{evenCount}(70,27)=29,\qquad T^{70}(27)=1.
\end{equation}
Сорок один ход двигателя против всего двадцати девяти рывков каната; канат в сумме отстаёт на
двенадцать и ни разу не вырывается вперёд из старта. \green{}
\end{theorem}

\emph{Идея доказательства.} Ядро закрывает префикс до $k=70$ прямым вычислением
(\code{counts\_le\_70\_at\_27}: $\mathrm{even}\le\mathrm{odd}$ в каждом префиксном окне;
\code{counts\_at\_70\_at\_27}: ровно $41$ против $29$ и финиш в $1$). Хвост падает в цикл
$1\to 2\to 1$, где шаги идут строго парами; лемма цикла \code{cycle\_counts} через аддитивность
счётчиков (\code{oddCount\_add}, \code{evenCount\_add}) замораживает дефицит $-12$ навсегда, так что ни
одно окно от старта его не перевернёт.

\begin{theorem}[\code{not\_ropeCountingLaw\_27}]\label{thm:not_ropeCountingLaw_27}
Закон каната ложен для $n=27$: $\neg\,\mathrm{RopeCountingLaw}(27)$. Из $j=0$ (значение $27>2$) ни одно
окно не даёт перевеса каната. \green{}
\end{theorem}

\begin{theorem}[\code{ropeLaw\_universal\_refuted}]\label{thm:ropeLaw_universal_refuted}
Универсальная форма закона ложна: $\neg\,\forall n\,(1\le n\Rightarrow\mathrm{RopeCountingLaw}(n))$,
свидетель $n=27$ (\cref{thm:not_ropeCountingLaw_27}). Граница-декрет на этом законе невозможна. \green{}
\end{theorem}

Красивая метафора «топливо кончится --- двигатель стянется к старту» верна по \emph{исходу} ($27$
всё-таки достигает $1$), но не по той сильной \emph{пошаговой} причине, которую формулировал закон. И
$27$ не одинок: харнесс находит целое семейство карабкающихся нарушителей ($27,31,41,47,54,55,\dots$).

\subsection{Неубывающая орбита есть вечный двигатель}

\begin{theorem}[\code{nonDescendingOrbit\_is\_perpetual\_engine}]\label{thm:nonDescendingOrbit_is_perpetual_engine}
Для $n\ge 3$ неубывающая орбита ($n\le T^{k}(n)$ для всех $k$) не проигрывает канату ни одного окна
($\forall k,\ \mathrm{evenCount}(k,n)\le\mathrm{oddCount}(k,n)$) и никогда не останавливается
($\forall K,\ T^{K}(n)\neq 1$). \green{}
\end{theorem}

\emph{Идея доказательства.} Неубывание вынуждает бюджет окна \cref{eq:window-budget} никогда не давать
чётному счётчику превзойти нечётный; бесконечная подпитка без единого проигранного окна и есть вечный
двигатель, а \cref{thm:window_descends} показывает, что до $1$ он не дойдёт.

\begin{theorem}[\code{nonHalting\_carries\_perpetual\_engine}]\label{thm:nonHalting_carries_perpetual_engine}
Если орбита никогда не достигает $1$, то её хвост от позиции глобального минимума --- неубывающая,
никогда не останавливающаяся орбита, предъявленная явно, а не через противоречие. \green{}
\end{theorem}

Натуральный ряд вполне упорядочен, так что у всякой орбиты есть глобальный минимум
(\code{exists\_min\_position}, стандартная тройка); хвост от него никогда не опускается ниже старта.
Прочитаем вслух: \textbf{всякий контрпример Коллатца несёт в себе вечный двигатель.} Гипотеза не
\emph{похожа} на утверждение «вечного двигателя нет» --- для карты $T$ она им \emph{является}, дословно:
орбита достигает $1$ тогда и только тогда, когда у неё нет вечного хвоста.

\subsection{Проверка, а не вывод}

\emph{Узнать причину изнутри} --- вывести универсальный закон каната внутренним доказательством ---
значило бы извлечь бесконечный факт о всей числовой прямой, не выходя за её пределы. Формализуем
попытку как самообоснование, пересекающее собственную границу.

\begin{definition}[\code{InternalisedCollatzGround}]\label{def:internalisedCollatzGround}
Структура с двумя полями: \code{ground} --- сам закон $\forall n\ge 1,\ \mathrm{RopeCountingLaw}(n)$, и
\code{beyondOwnHorizon} --- свидетельство того, что он выведен пересечением границы,
$\neg\,\forall n\ge 1,\ \mathrm{RopeCountingLaw}(n)$. Сокращение
\code{InternalKnowledgeOfCollatzCause} отождествляет внутреннее знание причины с этим самообоснованием.
\end{definition}

\begin{theorem}[\code{collatzCause\_unknowable}]\label{thm:collatzCause_unknowable}
Внутреннее знание первопричины Коллатца невозможно:
$\neg\,\mathrm{InternalKnowledgeOfCollatzCause}$. \green{} (вообще без аксиом)
\end{theorem}

\begin{proof}
Два поля \cref{def:internalisedCollatzGround} противоречат друг другу напрямую:
$\texttt{beyondOwnHorizon}\ \texttt{ground}:\mathrm{False}$.
\end{proof}

После падения декрета это держится a fortiori и \emph{по содержанию}: поле \code{ground} теперь
опровергнуто напрямую (\cref{thm:ropeLaw_universal_refuted}), а не только самоуничтожением формы. Это
тот же запрет, что не пускает вечный двигатель на $\N$ (\code{no\_infinite\_descent}) и, прочитанный по
Ландауэру, запрещает извлечь причину изнутри --- знание физично.

\begin{theorem}[\code{collatz\_verification\_not\_derivation}]\label{thm:collatz_verification_not_derivation}
Конъюнкция трёх звеньев: (1) бесконечное внутреннее извлечение закона каната невозможно,
$\neg\,\mathrm{InternalKnowledgeOfCollatzCause}$; (2) единственный внутренний след контрпримера ---
вечный двигатель (\code{nonHalting\_carries\_perpetual\_engine}); (3) значит, решение либо непознаваемо
изнутри, либо доступно лишь \emph{проверкой} орбиты. \green{}
\end{theorem}

Харнессы по $n\in[2,200000]$ со стражем расходимости на $10^5$ шагах --- это и есть «как далеко»:
никакое конечное окно не может \emph{доказать} универсальный закон --- оно может лишь не найти
опровержения на своём горизонте. А в пределе универсальный закон всё равно ложен --- не потому, что
орбита расходится, а потому что закон требовал большего, чем верно; проверка сходимости и проверка
закона --- не одно и то же.

\subsection{Декрет взят, затем снят; открытый статус}

Закон каната был принят \emph{декретом} --- четвёртым полем
$\code{collatzBoundary}:\forall n\ge 1,\ \mathrm{RopeCountingLaw}(n)$ единственной аксиомы \axiom{}, при
благоприятном знаке эвристики и в отсутствие тогда известного кованого опровержения. Его цена
раскрывалась сразу: доказана лишь стрелка «граница $\Rightarrow$ гипотеза», так что декрет мог
\emph{переплачивать}. Рядом были поставлены две растяжки --- намеренные детекторы взрыва вида
«опровержение закона $\Rightarrow \mathrm{False}$ ровно здесь». Когда \cref{thm:ropeLaw_universal_refuted}
был выкован, растяжка сработала: держать поле значило сделать аксиому противоречивой. Поле и его
растяжки покинули аксиому вместе, и границ снова три (твины, Риман, Навье--Стокс). Таинт не растёт;
после снятия он вернулся к $47$. Так трилемма Коллатца закрывается ровно как у Янг--Миллса --- кованым
опровержением, а не доказательством.

\begin{theorem}[\code{collatz\_open\_status}]\label{thm:collatz_open_status}
Эпистемический статус Коллатца после падения декрета (зеркало
\code{pnp\_locked\_behind\_engine\_status}, но уже \emph{без} декрет-конъюнкта) --- конъюнкция четырёх:
$\neg\,\forall n\ge 1,\ \mathrm{RopeCountingLaw}(n)$ (универсальный закон опровергнут) $\wedge$
$\neg\,\mathrm{InternalKnowledgeOfCollatzCause}$ (внутреннее знание невозможно) $\wedge$
$\bigl(\forall n\ge 1,\ \mathrm{RopeCountingLaw}(n)\Rightarrow\exists K,\ \mathrm{iter}\,K\,n=1\bigr)$
(per-$n$ закон влечёт остановку, условно) $\wedge$
$\bigl(\forall n,\ (\forall K,\ \mathrm{iter}\,K\,n\neq 1)\Rightarrow \exists j,\
\mathrm{NonDescendingOrbit}(\mathrm{iter}\,j\,n)\wedge\forall K,\ \mathrm{iter}\,K\,(\mathrm{iter}\,j\,n)\neq 1\bigr)$
(опровержение строило бы вечный двигатель). \green{}
\end{theorem}

Это не гёделевская независимость --- мы не утверждаем, что Коллатц недоказуем в арифметике Пеано.
Доказано другое, внутри нашей системы и строго: оба внутренних пути (опровержение и самообоснование)
упираются в один и тот же запрещённый объект, а декретный третий путь машинно закрыт. Сама гипотеза
остаётся открытой.

\begin{conjecture}[Коллатц]\label{conj:collatz}
Для всякого $n\ge 1$ найдётся $K$ с $T^{K}(n)=1$; равносильно, у всякой орбиты нет вечного хвоста.
\red{} Открыта и падением декрета не тронута: опровергнутый закон был строго сильнее гипотезы, так что
его падение говорит лишь, что \emph{таким} декретом Коллатца не закрыть.
\end{conjecture}

\subsection{Почему Коллатц стоит особняком}
Декрет пал именно здесь, и --- машинно проверено по всей программе --- больше нигде пасть не может. Атака того рода, что
свалила канат-закон, требует двух совпадений. Декретируемый закон должен быть \emph{строго сильнее} гипотезы, которую
охраняет --- доказана лишь стрелка (закон $\Rightarrow$ гипотеза), обратная ложна, --- так что контрпример убивает закон,
но оставляет гипотезу открытой; и его предикат должен быть \emph{разрешим} над $\mathbb{N}$, чтобы один малый свидетель
закрылся через \code{decide}. У канат-закона есть и то и другое: лестница \code{reaches\_one\_of\_countingLaw} идёт в одну
сторону, пустой спутник \code{valueLaw\_iff\_reaches\_one} вбирает эквивалентность сходимости, а $n=27$ --- который
\emph{сходится}, но закон нарушает --- опровергает его вычислением ядра.

Всякая другая граница проваливает хотя бы одно условие, и провал сам есть теорема. \emph{Взятые} декреты либо заперты
эквивалентностью, либо защищены барьером. Риман: \code{manifestationLaw\_iff\_RH\_of\_boundary} делает закон
$\Leftrightarrow$ RH, так что опровергнуть его --- \emph{значит} предъявить внекритический ноль, то есть опровергнуть RH,
над аналитической, не-$\mathbb{N}$ областью без \code{decide}; Янг--Миллс и Ходж --- то же
(\code{quantizationLaw\_iff\_massGap}, \code{descentLaw\_iff\_hodgeProperty}). Узел близнецов равносилен по силе своей
гипотезе (\code{quarantine\_inconsistent\_if\_node\_refuted}), и --- в отличие от безусловной у Коллатца
\code{nonHalting\_carries\_perpetual\_engine} --- его стрелка «опровержение строит двигатель» требует той самой
стабильности, что даёт декрет (\code{twinRefutation\_in\_stableUniverse\_builds\_engine}); барьер несовместимости
(\code{higherEnergyIncompatibility\_main}) запрещает внутреннее опровержение вовсе. У Навье--Стокса $\Leftrightarrow$-зеркала
нет совсем: единственный свидетель --- гладкое бессиловое решение на $\mathbb{R}^3$ с приростом энергии, не проще открытой
стационарной проблемы Лиувилля.

\emph{Отложенные} границы не мягче. Мерсенн и Ферма доказанно эквивалентны своим гипотезам под границей
(\code{mersenneManifestationLaw\_iff\_unbounded\_of\_boundary}, \code{fermatManifestationLaw\_iff\_unbounded\_of\_boundary});
их придержали по \emph{знаку} эвристики, а не потому что атакуемы. Разрешимые фронты зоопарка --- Гольдбах, Лежандр,
совершенные числа --- несут свидетель под \code{decide} (сильнейшая непредъявимость), но опять эквивалентный универсал, так
что опровержение было бы полным опровержением гипотезы.

Мораль заостряет урок раздела. Декрет может пасть честно лишь там, где он сверх-притязает на \emph{опровержимую} силу за
пределами своей гипотезы. То, что дисциплина принимает границу лишь когда её опровергающий свидетель непредъявим, либо её
закон эквивалентен гипотезе, --- ровно и делает всякий выживший декрет неатакуемым, и ровно оно поймало Коллатца.

Классическая формулировка принадлежит Коллатцу и обозрена Лагариасом~\cite{lagarias1985,lagarias2010}; сильнейший известный результат в её сторону --- теорема Тао о том, что почти все орбиты достигают почти ограниченных значений~\cite{tao2022}.
Урок раздела методологический: декрет, который пал честно --- с раскрытой ценой, со взведённой
растяжкой, с падением, зафиксированным теоремой, --- стоит больше декрета, который стоит
неисследованным. Машина, проверяющая любую гипотезу, сама есть тот объект, который она не может
превзойти.
